\section{Yöntem}
\label{sec:method}

\subsection{Genel mimari}
Önerilen model dört aşamadan oluşur: (1)~paralel ViT omurgalarıyla öznitelik
çıkarımı, (2)~dal başına izdüşüm, (3)~füzyon ve (4)~MLP sınıflandırıcı.
Şekil~\ref{fig:arch} bu akışı şematik olarak gösterir.

\begin{figure}[H]
  \centering
  \begin{tikzpicture}[
      >={Stealth[length=2.2mm]},
      box/.style={draw, rounded corners, align=center, minimum height=0.95cm,
                  minimum width=2.1cm, font=\small, inner sep=2pt},
      bb/.style={box, fill=blue!8},
      pr/.style={box, fill=green!8, minimum width=1.9cm},
      fu/.style={box, fill=orange!14, minimum width=1.9cm},
      cl/.style={box, fill=red!8},
      io/.style={align=center, font=\small},
      hdr/.style={font=\footnotesize\itshape, text=black!65, align=center},
      dim/.style={font=\scriptsize, text=black!70, midway, above},
      x=1cm, y=1cm]
    % nodes (column x, row y)
    \node[io] (img) at (0,0) {Görüntü\\$224{\times}224$};
    \node[bb] (vit)  at (3,1.7)  {ViT-B/16\\768-d};
    \node[bb] (swin) at (3,0)    {Swin-T\\768-d};
    \node[bb] (beit) at (3,-1.7) {BEiT-B/16\\768-d};
    \node[pr] (pv) at (5.9,1.7)  {İzdüşüm\\512-d};
    \node[pr] (ps) at (5.9,0)    {İzdüşüm\\512-d};
    \node[pr] (pb) at (5.9,-1.7) {İzdüşüm\\512-d};
    \node[fu] (fus) at (8.7,0) {Ağırlıklı\\füzyon\\512-d};
    \node[cl] (mlp) at (11.2,0) {MLP\\(256, dr.\ 0,3)};
    \node[box, fill=gray!12] (out) at (13.7,0) {Softmax\\23 sınıf};
    % stage headers
    \node[hdr] at (3,2.95)   {Öznitelik çıkarımı\\(omurga)};
    \node[hdr] at (5.9,2.95) {İzdüşüm};
    \node[hdr] at (8.7,2.95) {Füzyon};
    \node[hdr] at (12.45,2.95) {Sınıflandırma};
    % input fan-out (same image to 3 backbones)
    \draw[->] (img) -- (vit);
    \draw[->] (img) -- (swin);
    \draw[->] (img) -- (beit);
    % backbone -> projection (768-d)
    \draw[->] (vit) -- (pv);
    \draw[->] (swin) -- (ps) node[dim]{768-d};
    \draw[->] (beit) -- (pb);
    % projection -> fusion fan-in (512-d)
    \draw[->] (pv) -- (fus);
    \draw[->] (ps) -- (fus) node[dim]{512-d};
    \draw[->] (pb) -- (fus);
    % fusion -> MLP -> output
    \draw[->] (fus) -- (mlp) node[dim]{512-d};
    \draw[->] (mlp) -- (out) node[dim]{23};
  \end{tikzpicture}
  \caption{Üçlü ViT füzyon mimarisi. Aynı $224{\times}224$ görüntü üç omurgaya
  verilir; her omurga \texttt{timm} üzerinden \texttt{num\_classes=0} ile yüklenip
  768 boyutlu havuzlanmış öznitelik üretir. Her dal 512 boyuta izdüşürülür,
  ağırlıklı füzyonla birleştirilir ve tek bir MLP (gizli 256, dropout 0,3) ile 23
  sınıfa sınıflandırılır. Omurgalar dondurulmuş öznitelik çıkarımı veya son
  bloklarının ince ayarı ile kullanılır.}
  \label{fig:arch}
\end{figure}

\subsection{Omurgalar ve öznitelik çıkarımı}
Üç omurga, farklı tümevarımsal önyargılara ve ön-eğitim rejimlerine sahip olacak
biçimde seçilmiştir. ViT~\cite{dosovitskiy2021vit} görüntüyü yamalara bölüp standart
bir dönüştürücü kodlayıcıya verir; evrişimsel yerellik ve öteleme-değişmezlik
önyargısı taşımadığından küresel öz-dikkatle uzun-erimli ilişkileri modeller, fakat
güçlü temsil için büyük ölçekli ön-eğitime ihtiyaç duyar.
Swin~\cite{liu2021swin} yamaları derin katmanlarda birleştirip kaydırmalı-pencere
öz-dikkati kullanarak \emph{hiyerarşik} ve yerel-çok ölçekli temsiller üretir; bu,
evrişime benzer bir yerellik önyargısını geri kazandırır ve doku ile sınır gibi
yerel ipuçlarına duyarlıdır. BEiT~\cite{bao2022beit} ise maskeli görüntü
modellemesiyle (yamaların bir bölümünü maskeleyip görsel belirteçleri geri kazanma)
\emph{öz-denetimli} ön-eğitilir ve denetimli ViT'ten farklı bir temsil öğrenir. Bu
üç ayrı önyargı, füzyon için tamamlayıcı sinyaller sağlar.

Her omurga \texttt{timm} kütüphanesinden ImageNet üzerinde ön-eğitilmiş ağırlıklarla
yüklenir. Sınıflandırma başlığı \texttt{num\_classes=0} ile kaldırılır ve mimariye
uygun yerel havuzlanmış öznitelik vektörü ($\mathbb{R}^{768}$) elde edilir. Bu
yaklaşım hem CLS-belirteçli (ViT, BEiT) hem de CLS-belirteçsiz (Swin, küresel
ortalama havuzlama) mimarilerde doğru çalışır.

\subsection{İzdüşüm, füzyon ve sınıflandırıcı}
Bu aşamada her omurganın 768 boyutlu özniteliği ortak bir uzaya indirgenir, üç dal
tek bir vektörde birleştirilir ve sonuç sabit bir sınıflandırıcıya verilir.

\subsubsection{Dal izdüşümü}
Her dalın özniteliği, ardışık \emph{Doğrusal $\to$ LayerNorm $\to$ GELU}
katmanıyla 512 boyutlu ortak bir uzaya izdüşürülür.

\subsubsection{Füzyon yöntemleri}
Üç füzyon yöntemi uygulanmıştır:
\begin{itemize}
  \item \textbf{Birleştirme:} dal öznitelikleri uç uca eklenir ($N{\times}512$).
  \item \textbf{Ağırlıklı:} öğrenilen skaler ağırlıklarla normalize toplam (512).
  \item \textbf{GMU:} her dala öğrenilen \emph{öğe bazında} (per-feature) kapılarla
        ağırlık veren kapılı füzyon~\cite{arevalo2017gmu}. İki dal için bu, özgün
        makalenin $h = z \odot h_v + (1-z)\odot h_t$ biçimine \emph{tam olarak}
        indirgenir.
\end{itemize}
İzdüşürülmüş dal öznitelikleri $h_i \in \mathbb{R}^{512}$ ($i = 1,\dots,N$) için
füzyon işlemleri şöyle tanımlanır:
\begin{align}
  \text{Birleştirme:}\quad & f = [\,h_1;\, h_2;\, \dots;\, h_N\,] \in \mathbb{R}^{N\cdot 512}, \\
  \text{Ağırlıklı:}\quad & \alpha = \mathrm{softmax}(w),\qquad f = \sum_{i=1}^{N} \alpha_i\, h_i, \\
  \text{GMU:}\quad & z = \mathrm{softmax}_{\text{dal}}\!\big(W_z\,[h_1;\dots;h_N]\big),\qquad
                     f = \sum_{i=1}^{N} z_i \odot h_i.
\end{align}
Ağırlıklı füzyonda $w \in \mathbb{R}^{N}$ öğrenilen skaler ağırlıklardır ($\alpha_i$
skalerdir); GMU'da $z_i \in \mathbb{R}^{512}$ öğe-bazında kapılardır ve dal ekseninde
normalize edilir, dolayısıyla her öznitelik boyutu için ayrı ağırlıklandırma yapılır.
İki dal durumunda GMU, $f = z \odot h_1 + (1-z) \odot h_2$ biçimine indirgenir.

\subsubsection{Sınıflandırıcı}
Füzyon çıktısı tek bir MLP'ye ($256$ gizli birim, dropout $0{,}3$) verilir.
Tüm yapılandırmalarda sınıflandırıcı sabittir.

\subsection{Transfer öğrenme rejimleri}
İki rejim karşılaştırılmıştır: (i)~\emph{dondurulmuş} öznitelik çıkarımı (omurga
ağırlıkları sabit, yalnızca izdüşüm + füzyon + MLP eğitilir) ve (ii)~\emph{ince
ayar}: ViT/BEiT'te son üç dönüştürücü bloğu, Swin'de son aşama açılır; ayrıca son
normalizasyon katmanı eğitilir. Yalnız son blokların açılması bilinçli bir tercihtir:
küçük ve dengesiz bir veri kümesinde tüm ağı açmak, ImageNet ön-eğitiminde edinilen
genel öznitelikleri bozabilecek felakete-varan unutmaya (catastrophic forgetting)
yol açar; buna karşılık son bloklar göreve/alana özgü olduğundan uç (endoskopi)
alanına uyarlanmaya en uygun katmanlardır. Aynı gerekçeyle ince ayar, katman bazında
azalan öğrenme oranı (LLRD) ile uygulanır~\cite{howard2018ulmfit}: çıkışa yakın
katmanlar daha yüksek, gövdeye yakın (daha genel) katmanlar daha düşük öğrenme
oranıyla güncellenir, böylece genel temsiller korunurken üst katmanlar alana uyarlanır.

\subsection{Yorumlanabilirlik ve istatistiksel analiz}
Nihai model için (a)~CLS-belirteçli ViT-B ve BEiT-B dallarında \emph{dikkat
yayılımı}~\cite{abnar2020rollout}: ham dikkat, derin katmanlarda belirteç kimliğini
yitirdiğinden tek başına güvenilir bir açıklama vermez; kalıntı bağlantılarını
hesaba katmak için her katmanda $A = 0{,}5\,W_\text{att} + 0{,}5\,I$ alınır (birim
matris kalıntı yolunu temsil eder), katmanlar boyunca özyinelemeli çarpılır ve CLS
satırı $14{\times}14$ ızgaraya eşlenir; (b)~üç dalda
\emph{Grad-CAM++}~\cite{chattopadhyay2018gradcampp};
(c)~512 boyutlu füzyon öznitelikleri üzerinde \emph{UMAP}~\cite{mcinnes2018umap}
görselleştirmesi; (d)~eşli \emph{McNemar} testi~\cite{dietterich1998mcnemar} ve
(e)~son-işlem \emph{logit düzeltmesi}~\cite{menon2020logitadjust} uygulanmıştır.
Tüm bu çözümlemeler, eğitilmiş kontrol noktalarından \emph{yalnızca çıkarım}
yoluyla üretilmiş ve nihai modeli değiştirmemiştir.
